\section{Discussion}

In this paper, we presented a scalable fine-tuning method to improve multi-turn decision making abilities of LLMs. Moreover, we showed that the strategies learned by the LLM from our method can generalize zero-shot to unseen tasks. There are a few limitations to our approach. Firstly, we use rejection sampling on self-generated data to teach the model better behaviors. In order to get good performance, the starting model need to exhibit good behavior within a reasonable generation budget, so \ours{} would perform worse in the absence of a good base model. Next, we use offline preference tuning algorithms to train our models due to the lack of computational resources. A possible future direction for our work is to run online RL on diverse tasks instead: due to its recent success in other domains~\citep{deepseekai2025deepseekr1incentivizingreasoningcapability}, we expect it will give a larger improvement in LLMs' in-context RL capabilities. Our environments, despite being designed with the help of GPT-4o-mini, required a lot of human effort for implementation. A new axis of improvement can be training an LLM to scalably generate suitable tasks that can then be used to train the agent. Finally, the performance of our curriculum learning algorithm heavily depends on the quality of the task group clusters which is not ideal, and one can study possible improvements of this algorithm. We leave these directions for future work.

\section*{Impact Statement}

Our work can be used to train large language models that have better strategic exploration and decision making capabilities, which can have potential impact in the real world if agentic systems become wide spread. Our experiments are conducted in relatively simple and controlled environments and it is an open question what kind of impacts truly agentic systems will have on society. Other than that, this paper presents work whose goal is to advance the field of Machine Learning. There are many potential overall societal consequences of our work, none of which we feel must be specifically highlighted here.

\section*{Reproducibility Statement}

We provide sufficient details about our implementation, hyperparameters, environment design and dataset construction in the main paper and the appendix to effectively reproduce the results in this paper. Our code, training dataset and models can be found via the project website: \url{https://paprika-llm.github.io/}

\subsection*{Acknowledgement}

This work was supported in part by the U.S. Army Futures Command under Contract No. W519TC-23-C-0030. Moreover, it has greatly benefited from using the Delta advanced computing and data resource supported by the National Science Foundation (OAC 2005572) and the State of Illinois, as part of ACCESS-approved compute grants~\citep{access_compute}. Subsequent larger scale experiments on Gemma-3-12B-IT models were run using Bridges-2~\citep{psc_computing} at Pittsburgh Supercomputing Center through ACCESS  allocation CIS240901 from the Advanced Cyberinfrastructure Coordination Ecosystem: Services \& Support (ACCESS) program, which is supported by National Science Foundation grants \#2138259, \#2138286, \#2138307, \#2137603, and \#2138296. The authors thank Brandon Pusateri, Jillian Lehosky and Greg Bauer from ACCESS Support Staff for their incredible help at approving supplements and renewals for ACCESS compute grants throughout this project. Moreover, the work would not have finished so quickly without the help of Brett Bode from NCSA Delta Support Staff, who provided the authors critical help about properly utilizing the Delta cluster.  FT and YJ gratefully acknowledge Samuel Sokota, Daman Arora, Andrea Zanette, Yuda Song, Gaurav Ghosal, Yutong He, So Yeon Min, Kevin Li, Wen-Tse Chen, Xintong Duan and other members of Russ, Auton, Locus and AIRe lab for feedback received on an earlier versions of this work. FT greatly benefited from his discussions with Prof. Aviral Kumar and his lab's computational resources. YJ gratefully acknowledges the support of the Google PhD Fellowship.
